\begin{frame}{RDF — the triple}
  \framesubtitle{Everything is \texttt{subject} \texttt{predicate} \texttt{object}}

  \vspace{0.3em}
  \centering
  \begin{tikzpicture}[
    box/.style={draw, rounded corners=3pt, fill=newlightblue,
                text=newgray, font=\small, inner sep=6pt, align=left,
                text width=4.2cm},
    lit/.style={draw, rounded corners=3pt, fill=orange!15,
                text=newgray, font=\small, inner sep=6pt, align=left,
                text width=2.6cm}
  ]
    \node[box] (s) at (-1.2,0)
          {\texttt{<http://example.org/cim/Substation-42>}};
    \node[box] (t1) at (6.3,1.3)
          {\texttt{cim:Substation}};
    \node[lit] (t2) at (6.3,-1.3)
          {\texttt{"Main Street Substation"}};

    \node[font=\scriptsize\itshape, text=newmidblue, above left=1pt and -3pt of s]
         {subject};

    \draw[-latex, thick, newdarkblue]
      (s.east) to[bend left=12]
        node[pos=0.5, above, font=\scriptsize\itshape, text=newmidblue]
        {\texttt{rdf:type}}
      (t1.west);
    \draw[-latex, thick, newdarkblue]
      (s.east) to[bend right=8]
        node[pos=0.5, below=10pt, font=\scriptsize\itshape, text=newmidblue]
        {\texttt{cim:IdentifiedObject.name}}
      (t2.west);
  \end{tikzpicture}

  \vspace{0.5em}
  \footnotesize
  Same \texttt{subject} in both triples.
  The first object is a \emph{resource} (\texttt{cim:Substation});
  the second is a \emph{literal} (a string).
\end{frame}
